% Préambule commun aux cours IPR Scholar.
\documentclass[11pt,a4paper,oneside]{article}
\usepackage[T1]{fontenc}
\usepackage[utf8]{inputenc}
\usepackage[french]{babel}
\usepackage{lmodern}
\usepackage{microtype}
\usepackage[a4paper,margin=2.35cm]{geometry}
\usepackage{amsmath,amssymb,mathtools}
\usepackage{booktabs,longtable,array}
\usepackage{xcolor}
\usepackage{hyperref}
\usepackage{tikz}
\usetikzlibrary{arrows.meta,positioning,shapes.geometric,calc}
\hypersetup{colorlinks=true,linkcolor=blue!55!black,citecolor=blue!55!black,urlcolor=blue!55!black}
\definecolor{iprblue}{HTML}{1F4E79}
\definecolor{iprgreen}{HTML}{2D6A4F}
\newcommand{\IPR}{\textsc{ipr}}
\newcommand{\coursmeta}[2]{\author{Projet \IPR{} -- documentation Scholar}\date{#1}\maketitle\begin{center}\small #2\end{center}\vspace{0.5cm}\tableofcontents\clearpage}
\newcommand{\resultat}[1]{\paragraph{À retenir.} #1}

\title{Cours 6 -- Correction de voile : diffusion atmosphérique et Dark Channel Prior}
\begin{document}
\coursmeta{3 octobre 2026}{Modèle physique simplifié et estimation statistique employés par le réglage Voile.}

\section{Modèle physique}
Dans un milieu diffusant, la lumière observée est la combinaison de la radiance de scène atténuée et de la lumière voilante de l'atmosphère :
\begin{equation}
I(x)=J(x)t(x)+A(1-t(x)).
\end{equation}
$I$ est le pixel observé, $J$ la scène recherchée, $A$ l'éclairement atmosphérique et $t\in[0,1]$ la transmission. Si $t$ est faible, les contrastes diminuent et les couleurs convergent vers $A$.

Le \emph{Dark Channel Prior} (DCP) constate que, hors ciel, une petite fenêtre de scène naturelle contient souvent au moins un canal sombre. Le canal noir est $D(x)=\min_{y\in\Omega(x)}\min_c I_c(y)$. Une érosion morphologique approxime cette recherche locale~\cite{he2009,he2011}.

\section{Application dans \IPR{}}
Le moteur calcule le minimum RGB, l'érode avec un noyau carré impair, estime $A$ au percentile $99$ de sous-échantillons lumineux, puis forme une transmission
\begin{equation}
t=\operatorname{clamp}(1-0{,}85\,a\,D/A,\;0{,}12,\;2).
\end{equation}
L'inversion est $J=(I-A)/t+A$. Un réglage négatif ajoute du voile; un réglage positif en retire. Les canaux finaux sont bornés à zéro, la compression de gamut étant reportée à l'export.

\section{Exemple}
Dans un paysage brumeux, les zones éloignées ont un minimum RGB élevé et une transmission faible : le rendu augmente leur contraste. Dans une scène de studio blanche, l'hypothèse DCP est faible; le même réglage peut créer une couleur artificielle. C'est pourquoi ce contrôle reste global et modéré.

\section{Pourquoi les constantes sont bornées}
Le maximum de 8192 échantillons borne le coût de l'estimation de percentile. Le noyau est limité à 15 pixels afin de ne pas transformer le voile en correction globale de grande échelle. Le plancher $t=0{,}12$ évite la division explosive qui amplifie bruit et artefacts. Le facteur $0{,}85$ est le compromis DCP classique entre correction et préservation d'une petite quantité de voile naturel.

\resultat{Le voile n'est pas une simple hausse de contraste : c'est l'inversion régularisée d'un modèle de diffusion, fiable surtout lorsque son hypothèse de scène est satisfaite.}
\bibliographystyle{plain}\bibliography{references}
\end{document}
